\documentclass{beamer}

\input{preamble}

\title[Beamer example]{Beamer template example}
\author{Author}
\institute{Institution}
\date{\today}

\showtodostrue

\begin{document}

\frame{\titlepage}

\begin{frame}
    \frametitle{Outline}
    \tableofcontents
\end{frame}

\section{Mathematical content}

\begin{frame}
    \frametitle{A theorem-like slide}

    \begin{theorem}[Sample theorem]
        Let \(X\) be a sample object. Then \(X\) satisfies a property visible in a
        Beamer theorem block.
    \end{theorem}

    \begin{proof}
        The proof is displayed using the template's mathematical font setup. For
        example,
        \[
            (a+b)^2 = a^2 + 2ab + b^2.
        \]
    \end{proof}

    \todo{Add more details to the proof.}
\end{frame}

\begin{frame}
    \frametitle{Lists and displayed mathematics}

    The template should handle ordinary mathematical slides with lists.

    \begin{itemize}
        \item First item, with inline mathematics \(f \colon X \to Y\).
        \item Second item, referring to the displayed equation below.
        \item Third item, with a citation \cite{har77}.
    \end{itemize}

    \[
        \chi(\mathscr{F}) =
        \sum_i (-1)^i h^i(X,\mathscr{F}).
    \]
\end{frame}

\section{A content-heavy slide}

\begin{frame}
    \frametitle{A denser slide}

    This slide is intentionally more text-heavy. It is meant to test line breaks,
    mathematical notation, lists, and the visual density of the theme.

    \begin{enumerate}
        \item The first point contains a short mathematical statement.
        \item The second point includes a displayed formula:
            \[
                0 \to \mathscr{F}' \to \mathscr{F} \to \mathscr{F}'' \to 0.
            \]
        \item The third point checks that the slide remains readable.
    \end{enumerate}
\end{frame}

\section{Bibliography}

\begin{frame}[allowframebreaks]
    \frametitle{References}
    \printbibliography
\end{frame}

\end{document}
